%=============================================================================
% APPENDIX R: EM-ψ BACK-REACTION COUPLING
% v3.0: New appendix from EM Coupling Bounds paper integration
%=============================================================================

\section{EM--$\psi$ Back-Reaction Coupling}\label{app:em-psi-coupling}

This appendix develops the framework for electromagnetic back-reaction on 
the scalar field $\psi$, introducing a single dimensionless parameter $\lambda$ 
that controls whether EM fields can source $\psi$ oscillations. We derive 
both ``accidental'' constraints from existing cavity stability and 
``intentional'' search protocols that could reach $|\lambda - 1| \sim 10^{-14}$.

%-----------------------------------------------------------------------------
\subsection{Physical Interpretation of $\lambda$}\label{sec:lambda-physics}

The parameter $\lambda$ toggles the EM--$\psi$ interaction:
\begin{itemize}
\item $\lambda = 1$: EM probes the optical metric $n = e^\psi$ but does not source $\psi$
\item $|\lambda - 1| \neq 0$: EM can pump $\psi$ modes (laboratory generation possible)
\end{itemize}

\paragraph{Intuitive picture.}
Think of $\psi$ as water and EM as a paddle:
\begin{itemize}
\item $\lambda = 1$: The paddle slides across without making waves
\item $|\lambda - 1| \neq 0$: The paddle makes waves; pump with the right rhythm and they grow
\end{itemize}

\paragraph{Relation to core postulates.}
The core \DFD postulates (Sec.~\ref{sec:core_idea}) specify how $\psi$ 
affects EM propagation ($n = e^\psi$, $c_1 = ce^{-\psi}$). The parameter 
$\lambda$ addresses the \textit{inverse} question: can EM fields actively 
modify $\psi$? This is a distinct physical degree of freedom not constrained 
by the forward propagation relations.

%-----------------------------------------------------------------------------
\subsection{Mode Equation and Pumping Channels}\label{sec:mode-equation}

\subsubsection{Single Lab-Mode Reduction}

Reduce the $\psi$ field to a single laboratory mode $q(t)$ with natural 
frequency $\Omega_\psi$ and damping $\gamma_\psi$:
\begin{equation}
\boxed{
\ddot{q} + 2\gamma_\psi\dot{q} + \Omega_\psi^2 q = \frac{(\lambda - 1)}{M_\psi}\int u(\mathbf{r})\,\Xi(\mathbf{r},t)\,d^3r + \alpha U(t)q
}
\label{eq:driven-mode}
\end{equation}
where:
\begin{itemize}
\item $u(\mathbf{r})$: normalized spatial profile of the $\psi$ mode
\item $M_\psi$: effective mass of the mode
\item $\Xi(\mathbf{r},t) \equiv -\frac{1}{2}e^{-2\psi_0}\left(B^2 - \frac{E^2}{c^2}\right)$: EM stress tensor trace
\item $U(t) = U_0[1 + m\cos(2\omega t)]$: stored EM energy with modulation depth $m$
\item $\alpha$: parametric coupling coefficient
\end{itemize}

The EM stress $\Xi$ carries a $2\omega$ component for a cavity driven at 
frequency $\omega$, providing two pumping channels.

\subsubsection{Channel 1: Driven Resonance ($2\omega = \Omega_\psi$)}

When twice the EM drive frequency matches the $\psi$-mode frequency, 
direct resonant driving occurs. The steady-state amplitude is:
\begin{equation}
|q|_{\rm res} \simeq \frac{|\lambda - 1||G|}{2M_\psi\Omega_\psi\gamma_\psi},
\label{eq:driven-amplitude}
\end{equation}
where the geometry overlap is:
\begin{equation}
G \equiv \int u(\mathbf{r})\,\hat{\Xi}_{2\omega}(\mathbf{r})\,d^3r,
\end{equation}
with $\hat{\Xi}_{2\omega}$ the $2\omega$ Fourier component of $\Xi$.

\subsubsection{Channel 2: Parametric Amplification ($2\omega \simeq 2\Omega_\psi$)}

The stiffness modulation from $U(t)$ creates parametric gain. The Mathieu 
gain parameter is:
\begin{equation}
h = (\lambda - 1)\frac{U_0}{M_\psi\Omega_\psi^2}Hm,
\label{eq:mathieu-h}
\end{equation}
with instability growth rate:
\begin{equation}
\Gamma \simeq \frac{1}{2}h\Omega_\psi - \gamma_\psi.
\end{equation}

The overlap $H$ is:
\begin{equation}
H = \frac{1}{U_0}\int u^2(\mathbf{r})\,w(\mathbf{r})\,d^3r, \quad 
w = \frac{\varepsilon_0}{4}E^2 + \frac{\mu_0}{4}H^2.
\label{eq:overlap-H}
\end{equation}

\paragraph{Instability threshold.}
Parametric instability occurs when $\Gamma > 0$:
\begin{equation}
|\lambda - 1|_{\rm min} = \frac{2\gamma_\psi}{\Omega_\psi}\frac{M_\psi\Omega_\psi^2}{U_0 H m}.
\label{eq:instability-threshold}
\end{equation}

%-----------------------------------------------------------------------------
\subsection{Geometry Transparency}\label{sec:geometry}

\subsubsection{When the Driven Overlap Cancels}

For a single, symmetric pillbox cavity driven in a pure eigenmode (TM$_{010}$ 
or TE$_{011}$), Bessel identities and time-averaged equipartition make:
\begin{equation}
\int \left(B^2 - \frac{E^2}{c^2}\right) d^3r \approx 0 \quad \Rightarrow \quad G \approx 0.
\end{equation}

The driven channel is \textit{geometrically transparent} for symmetric 
cavities in pure eigenmodes.

\subsubsection{How to Restore the Overlap}

Three methods restore $G \neq 0$:
\begin{enumerate}
\item \textbf{TE+TM superposition}: Co-phased modes with matched radii give 
$G = u(z_0)e^{-2\psi_0}\eta_\times U_0\cos\phi$, where $\eta_\times = \mathcal{O}(0.1$--$1)$.

\item \textbf{Asymmetric geometry}: Small irises or near-cutoff asymmetries 
break equipartition.

\item \textbf{Mode beating}: Two nearby modes at frequencies $\omega_1$, 
$\omega_2$ produce $2\omega = \omega_1 + \omega_2$ components.
\end{enumerate}

\subsubsection{Parametric Overlap: Robust Area-Ratio Law}

For a $\psi$-mode ``tube'' of height $L$ and cross-section $A_\psi$, with 
$N$ compact cavities of total aperture $A_{\rm cav,tot}$ placed at antinodes:
\begin{equation}
H \approx \frac{2}{L}\kappa_{\rm eff}\frac{A_{\rm cav,tot}}{A_\psi},
\label{eq:area-ratio-law}
\end{equation}
where $\kappa_{\rm eff} = \mathcal{O}(1)$ captures mode-shape details.

Combining with~\eqref{eq:instability-threshold} and using 
$M_\psi \simeq A_\psi L/(2\pi c_s)$ for a 1D standing mode:
\begin{equation}
\boxed{
|\lambda - 1|_{\rm min} = \frac{\pi\gamma_\psi}{c_s U_0 m}\frac{A_\psi^2}{\kappa_{\rm eff}A_{\rm cav,tot}}
}
\label{eq:design-law}
\end{equation}

%-----------------------------------------------------------------------------
\subsection{Constraints on $|\lambda - 1|$}\label{sec:lambda-bounds}

\subsubsection{Accidental Constraint from Cavity Stability}

The mere stability of existing high-Q cavities---the absence of observed 
parametric instability near twice the drive frequency---provides a 
conservative bound.

\paragraph{Conservative parameters.}
\begin{itemize}
\item Stored energy: $U_0 \sim 100$ kJ
\item Modulation depth: $m \sim 0.01$ (ambient amplitude/PLL dither)
\item Loss ratio: $\gamma_\psi/\Omega_\psi \sim 10^{-3}$
\item Tube area: $A_\psi \sim 0.8$ m$^2$
\item Cavity aperture: $A_{\rm cav,tot} \sim 3 \times 10^{-3}$ m$^2$
\item $\kappa_{\rm eff} \sim 1$, $c_s \leq c$
\end{itemize}

\paragraph{Result.}
Using Eq.~\eqref{eq:design-law}:
\begin{equation}
\boxed{|\lambda - 1| \lesssim 3 \times 10^{-5}}
\label{eq:accidental-bound}
\end{equation}

Any substantially larger coupling would have produced obvious parametric 
instability in normal cavity operation---and it has not.

\subsubsection{Intentional Search: Projected Reach}

With deliberate optimization using the same physics:
\begin{itemize}
\item $U_0 \to 1$ MJ (factor 10 increase)
\item $m \to 0.1$ (factor 10 increase)
\item Array apertures at all antinodes: $A_{\rm cav,tot} \to 3 \times 10^{-2}$ m$^2$ (factor 10)
\item Shrink tube area: $A_\psi \to 0.27$ m$^2$ (factor $\sim$3 reduction)
\item Maintain $\gamma_\psi/\Omega_\psi \sim 10^{-3}$
\end{itemize}

The design law~\eqref{eq:design-law} then gives:
\begin{equation}
\boxed{|\lambda - 1| \sim 10^{-14} \quad \text{(accessible reach)}}
\label{eq:intentional-reach}
\end{equation}

\begin{table}[h]
\caption{Accidental vs. intentional search parameters.}
\label{tab:lambda-parameters}
\begin{ruledtabular}
\begin{tabular}{lcc}
Parameter & Accidental & Intentional \\
\hline
Stored energy $U_0$ (J) & $10^5$ & $10^6$ \\
Modulation depth $m$ & 0.01 & 0.10 \\
Cavity aperture $A_{\rm cav,tot}$ (m$^2$) & $3 \times 10^{-3}$ & $3 \times 10^{-2}$ \\
Tube area $A_\psi$ (m$^2$) & 0.8 & 0.27 \\
Loss ratio $\gamma_\psi/\Omega_\psi$ & $10^{-3}$ & $10^{-3}$ \\
\hline
Projected $|\lambda - 1|_{\rm min}$ & $\lesssim 3 \times 10^{-5}$ & $\sim 10^{-14}$ \\
\end{tabular}
\end{ruledtabular}
\end{table}

%-----------------------------------------------------------------------------
\subsection{Why $\lambda \neq 1$ Has Not Been Detected}\label{sec:why-not-seen}

Three factors explain the null result in existing metrology:

\begin{enumerate}
\item \textbf{Pure eigenmodes suppress the driven channel.}
Symmetric cavities in pure modes have $G \approx 0$ by Bessel-function 
orthogonality and equipartition.

\item \textbf{Parametric pumping needs deliberate $2\omega$.}
Routine metrology avoids such tones and heavily filters them to suppress 
amplitude-modulation sidebands.

\item \textbf{$2\omega$ features treated as technical noise.}
Any residual $2\omega$ response is interpreted as technical AM sidebands 
and actively suppressed, not investigated as a potential signal.
\end{enumerate}

\paragraph{To detect $|\lambda - 1| \neq 0$:}
\begin{itemize}
\item Use TE+TM superposition (restores $G \neq 0$)
\item Deliberately apply $2\omega$ modulation
\item \textit{Preserve} (not suppress) $2\omega$ response
\item Monitor for resonant growth at $\Omega_\psi$
\end{itemize}

%-----------------------------------------------------------------------------
\subsection{Intentional Detection Protocol}\label{sec:detection-protocol}

\begin{tcolorbox}[colback=blue!5!white, colframe=blue!50!black, title=Intentional $\psi$-Pump Detection: Required Capabilities]
\begin{enumerate}
\item \textbf{High-Q resonator} ($Q \gtrsim 10^4$) with stored energy $U_0 \gtrsim 1$ MJ 
(pulsed acceptable)

\item \textbf{Phase-stable amplitude modulation} at $2\omega$ with depth $m \sim 0.1$ 
on stored energy

\item \textbf{Placement of cavity apertures at $\psi$ antinodes} (maximize $H$; 
use multiple irises)

\item \textbf{Phase-sensitive readout near $\Omega_\psi$}; preserve $2\omega$ tones 
(do not auto-suppress)

\item \textbf{Null sensitivity target}: $\Delta\psi \lesssim 10^{-14}$ or 
equivalently $|\lambda - 1| \lesssim 10^{-14}$
\end{enumerate}
\end{tcolorbox}

\paragraph{Orthogonal cross-check: Driven amplitude.}
With a TE+TM superposition ($\eta_\times \neq 0$, phase $\phi = 0$):
\begin{equation}
\Delta\psi \equiv u(z_0)|q|_{\rm res} \approx \frac{|\lambda - 1|\eta_\times U_0 c_s}{\pi A_\psi \gamma_\psi}.
\label{eq:driven-crosscheck}
\end{equation}
For $\eta_\times \sim 0.3$, $U_0 = 100$ kJ, $A_\psi = 0.8$ m$^2$, 
$\gamma_\psi = 0.03$ s$^{-1}$:
\begin{equation}
\Delta\psi \sim 1.2 \times 10^{-3}|\lambda - 1|,
\end{equation}
which crosses cavity-atom sensitivity (Sec.~\ref{sec:cavity}) in the 
$10^{-12}$--$10^{-15}$ range for $|\lambda - 1|$ in $10^{-9}$--$10^{-12}$.

%-----------------------------------------------------------------------------
\subsection{Relation to Core DFD Framework}\label{sec:lambda-relation}

\paragraph{Consistency with postulates.}
The parameter $\lambda$ does \textit{not} modify the core postulates:
\begin{itemize}
\item Refractive index: $n = e^\psi$ (unchanged)
\item One-way light speed: $c_1 = ce^{-\psi}$ (unchanged)
\item Matter acceleration: $a = (c^2/2)\nabla\psi$ (unchanged)
\item Field equation: Eq.~\eqref{eq:field_general} (unchanged for static/quasi-static)
\end{itemize}

The $\lambda$ parameter describes a \textit{dynamic} EM--$\psi$ interaction 
orthogonal to the static field relations. It affects how rapidly oscillating 
EM fields can pump $\psi$ modes, not how $\psi$ affects light propagation.

\paragraph{Default value.}
Without additional physics, $\lambda = 1$ (no back-reaction) is the natural 
default. Any $|\lambda - 1| \neq 0$ indicates additional EM--gravity coupling 
beyond metric propagation effects.

%-----------------------------------------------------------------------------
\subsection{Summary}\label{sec:lambda-summary}

\begin{enumerate}
\item The parameter $\lambda$ controls EM back-reaction on $\psi$: 
$\lambda = 1$ means EM probes but doesn't pump; $|\lambda - 1| \neq 0$ 
enables laboratory $\psi$-generation.

\item Existing cavity stability provides an \textbf{accidental bound}:
\begin{equation}
|\lambda - 1| \lesssim 3 \times 10^{-5}.
\end{equation}

\item Deliberate optimization enables an \textbf{intentional search} reaching:
\begin{equation}
|\lambda - 1| \sim 10^{-14}.
\end{equation}

\item The null detection so far is explained by geometry transparency and 
suppression of $2\omega$ components in standard metrology.

\item A dedicated search protocol with TE+TM superposition and preserved 
$2\omega$ response could either discover $\lambda \neq 1$ or constrain it 
below $10^{-14}$ using existing apparatus.
\end{enumerate}

\begin{tcolorbox}[colback=green!5!white, colframe=green!50!black, title=Key Result]
\textbf{We are not asking anyone to believe new physics; we are asking 
them to notice the parametric instability that is not there.}

Unoptimized cavities accidentally constrain $|\lambda - 1| \lesssim 3 \times 10^{-5}$. 
An intentional $2\omega$ modulation test using the same hardware pushes 
\textbf{ten orders of magnitude tighter}. A single afternoon's measurement 
could either discover $\lambda \neq 1$ or constrain it below $10^{-14}$.
\end{tcolorbox}

%-----------------------------------------------------------------------------
\subsection{Dual-Sector Extension: The $\kappa$ Parameter}\label{sec:kappa-parameter}
%-----------------------------------------------------------------------------

Beyond the $\lambda$ parameter controlling EM back-reaction, a second 
parameter $\kappa$ controls the \textit{differential} 
response of electric and magnetic sectors to $\psi$.

\paragraph{Status of the $\kappa$ parameter.}
The parameter $\kappa$ should not be viewed as a free phenomenological
constant at leading order. At tree level, the Gordon optical metric gives
$\kappa = 0$, i.e.\ no electric--magnetic constitutive split. Within the
gauge-emergence auxiliary-metric completion of DFD, however, a nonzero
split is induced, yielding the definite prediction
\begin{equation}\label{eq:kappa-derived}
\kappa = \alpha_{\rm eff} = \frac{\alpha}{n_2^2} = \frac{\alpha}{4}
\approx 1.82\times 10^{-3},
\end{equation}
where $n_2 = 2$ is the $\mathrm{SU}(2)$ frame stiffness associated with
the $(3,2,1)$ partition (Appendix~\ref{app:alpha-derivations};
see also Ref.~\cite{AlcockVacuumLoading2026}).
Existing cavity-stability bounds such as
$|\kappa| \lesssim 1$ should therefore be interpreted not as the primary
definition of $\kappa$, but as an independent experimental consistency
check on the derived prediction. The DFD hierarchy is:
tree-level Gordon sector $\kappa = 0$, gauge-emergence completion
$\kappa = \alpha/4$, and experiment tests consistency with that value.

\subsubsection{Constitutive Split Preserving $v_{\rm ph} = c/n$}

The vacuum permittivity and permeability can respond asymmetrically to $\psi$:
\begin{equation}\label{eq:constitutive-split}
\epsilon(\psi) = \epsilon_0\, n\, e^{+\kappa\psi}, \qquad
\mu(\psi) = \mu_0\, n\, e^{-\kappa\psi},
\end{equation}
where $n = e^\psi$ and $\kappa$ is the split parameter.

The product is preserved:
\begin{equation}
\epsilon(\psi)\mu(\psi) = \epsilon_0\mu_0\, n^2 \quad \Rightarrow \quad
v_{\rm ph} = \frac{1}{\sqrt{\epsilon\mu}} = \frac{c}{n}.
\end{equation}
Thus the optical metric phase speed is unchanged by the split.

\paragraph{Physical interpretation.}
\begin{itemize}
\item $\kappa = 0$: Electric and magnetic sectors respond identically to $\psi$ (symmetric case)
\item $\kappa \neq 0$: Sector-differential response; electric and magnetic energies couple differently
\end{itemize}

\subsubsection{The Unified Bracket}

With the split~\eqref{eq:constitutive-split}, a single bracket governs energy exchange, body force, and $\psi$ sourcing:
\begin{equation}\label{eq:unified-bracket}
\mathcal{B} \equiv \frac{B^2}{\mu} - \epsilon E^2.
\end{equation}

\paragraph{Energy exchange.}
The Poynting theorem acquires:
\begin{equation}
\partial_t u + \nabla\cdot\mathbf{S} = -\mathbf{J}\cdot\mathbf{E} - \frac{\kappa}{2}\dot{\psi}\,\mathcal{B}.
\end{equation}

\paragraph{Body force.}
Fields exert force on the medium:
\begin{equation}
\mathbf{f}_\psi = -\frac{\kappa}{2}\,\mathcal{B}\,\nabla\psi.
\end{equation}

\paragraph{$\psi$ sourcing.}
EM fields can source $\psi$:
\begin{equation}
\frac{\delta\mathcal{L}_\psi}{\delta\psi} = \mathcal{S}_{\rm mass} + \frac{\kappa}{2}\,\mathcal{B}.
\end{equation}

\subsubsection{Standing-Wave Energy Equality}

For a lossless, steady-state standing wave in a linear medium, the cycle-averaged \textit{integrated} energies are equal:
\begin{equation}
\int_V \overline{\epsilon E^2}\,dV = \int_V \overline{B^2/\mu}\,dV,
\end{equation}
so $\int_V \overline{\mathcal{B}}\,dV = 0$. The integrated bracket vanishes for ideal standing waves.

\paragraph{Local imbalance.}
Nonzero \textit{local} bracket arises at $\mathcal{O}(\theta^2)$ due to longitudinal fields in paraxial Gaussian modes. For a TEM$_{00}$ cavity mode with waist $w_0$:
\begin{equation}
\overline{\epsilon E^2} - \overline{B^2/\mu} \sim \theta^2\,\epsilon|E_0|^2, \qquad \theta = \frac{\lambda}{\pi w_0}.
\end{equation}
For $\lambda = 1064$ nm, $w_0 = 300\,\mu$m: $\theta^2 \simeq 1.3 \times 10^{-6}$.

\subsubsection{Experimental Tests of the $\kappa = \alpha/4$ Prediction}

\paragraph{Accidental bound from cavity stability.}
Absence of $2\omega$ parametric instabilities in extreme-Q resonators constrains unintended EM$\leftrightarrow\psi$ pumping. This provides headroom consistent with $|\kappa| \lesssim 1$, which is satisfied by $\kappa = \alpha/4 \approx 0.002$ by three orders of magnitude.

\paragraph{LPI residual as $\kappa$ test.}
After the constitutive-chain cancellation of Sec.~\ref{sec:cavity-proof}, the cavity--atom observable is a screened residual rather than an order-unity slope.  Nevertheless, the sector-resolved residual still depends on $\kappa$ via:
\begin{align}
\xi_{\rm LPI}^{\rm res}(\kappa)
  &= \text{(screened residual of)}\;
     1 - \alpha_L^{(M)} - \alpha_{\rm at}^{(S)}(\kappa),
     \notag\\
\alpha_{\rm at}^{(S)}(\kappa)
  &= K_\epsilon^{(S)}\kappa + \mathcal{O}(\kappa^2).
\end{align}
where $K_\epsilon^{(S)}$ is the atomic EM-energy sensitivity.  At leading order in the gauge-emergence completion, $\kappa$ is predicted to be $\alpha/4$; experiment serves to test this prediction and bound any higher-order or screening corrections.

\paragraph{Order-of-magnitude for $K_\epsilon^{(S)}$.}
Atomic optical transition energies scale with the effective Rydberg $R_\infty \propto 1/\epsilon^2$, giving:
\begin{equation}
\frac{\delta E}{E}\bigg|_{\rm gross} \simeq -2\,\frac{\delta\epsilon}{\epsilon} \quad \Rightarrow \quad K_\epsilon^{(S)} \sim \mathcal{O}(1\text{--}3).
\end{equation}
For Sr and Yb clock transitions, $K_\epsilon^{(S)}$ is plausibly order unity.

\subsubsection{Experimental Discrimination}

The prediction $\kappa = \alpha/4$ can be tested via:
\begin{enumerate}
\item \textbf{TE/TM polarization swaps:} Pure TE (magnetic dominant) vs pure TM (electric dominant) modes have opposite bracket signs.
\item \textbf{Dual-wavelength measurements:} $\kappa$ is wavelength-independent; dispersion effects are not.
\item \textbf{Multi-species clock comparisons:} Different atoms have different $K_\epsilon^{(S)}$ values.
\end{enumerate}

\begin{tcolorbox}[colback=blue!5!white, colframe=blue!50!black, title=Dual-Sector Extension Summary]
\textbf{The $\kappa$ parameter:}
\begin{itemize}[nosep]
\item Controls differential $\epsilon/\mu$ response to $\psi$ while preserving $v_{\rm ph} = c/n$
\item Unified bracket $\mathcal{B} = B^2/\mu - \epsilon E^2$ governs energy, force, and sourcing
\item \textbf{Predicted:} $\kappa = \alpha/4 \approx 1.82\times 10^{-3}$ from gauge-emergence completion
\item Consistent with cavity stability bound $|\kappa| \lesssim 1$; directly testable via sector-resolved LPI slope
\end{itemize}
\textbf{Falsification:} If TE/TM cavity comparisons show no $\psi$-dependent split at $10^{-5}$ precision, $\kappa \approx 0$ is confirmed and the gauge-emergence prediction $\kappa = \alpha/4$ is falsified.
\end{tcolorbox}
